\documentclass[11pt,letterpaper]{amsart}
\usepackage[T1]{fontenc}
\usepackage[utf8]{inputenc}
\usepackage{lmodern}
\usepackage{amsmath,amssymb,amsfonts}
\usepackage{geometry}
\usepackage{enumitem}
\usepackage{xcolor}
\usepackage[hidelinks]{hyperref}
\geometry{margin=1in}

\newcommand{\R}{\mathbb R}
\newcommand{\Len}{\mathrm{Len}}
\newcommand{\Roll}{\mathrm{Roll}}
\DeclareMathOperator{\Ker}{ker}
\DeclareMathOperator{\rank}{rank}

\newenvironment{checkpoint}
{\par\medskip\noindent\begin{quote}\textbf{Checkpoint.}\ }
{\end{quote}\par\medskip}

\newenvironment{warningbox}
{\par\medskip\noindent\begin{quote}\textbf{Advertencia.}\ }
{\end{quote}\par\medskip}

\newenvironment{taskbox}
{\par\medskip\noindent\begin{quote}\textbf{Tarea.}\ }
{\end{quote}\par\medskip}

\newenvironment{microbox}
{\par\medskip\noindent\begin{quote}\textbf{Micro-modulo.}\ }
{\end{quote}\par\medskip}

\title{Plan de desarrollo de \texttt{csdraw}\\
Software para medir longitud de diagramas cs de nudos}
\author{Guia concreta para Tomas}
\date{}

\begin{document}
\maketitle

\section{Objetivo del software}

El objetivo es construir software para medir, dibujar y luego optimizar la
longitud de diagramas cs de nudos.

La primera version debe tener dos tests obligatorios, en este orden:
\[
\boxed{
1.\ \texttt{stadium\_2.json}
}
\qquad
\boxed{
2.\ \texttt{figure8\_4\_1\_manual.json}
}
\]

El primer test, \texttt{stadium\_2.json}, no es un nudo. Es el test geometrico
minimo del motor cs. Verifica puntos, longitud, SVG, contactos, rolling,
kernel, gradiente reducido y descenso.

El segundo test, \texttt{figure8\_4\_1\_manual.json}, es el primer test real de
diagrama de nudo. Debe construirse a partir de un diagrama conocido del nudo
figura ocho \(4_1\), usando una componente declarada y cuatro cruces over/under
declarados manualmente.

\begin{warningbox}
No debemos confundir a Tomas. Primero se prueba el motor cs con
\texttt{stadium\_2.json}. Solo despues se prueba el diagrama de nudo
\texttt{figure8\_4\_1\_manual.json}. No usamos gauge. No usamos variables
\(\omega_\alpha\). No usamos \(L(c)\). No usamos codigo de Gauss en la primera
version.
\end{warningbox}

\begin{warningbox}
En la primera version, \(4_1\) se usa para validar estructura de diagrama de
nudo, no para certificar estacionariedad ni descenso. Los modulos
\texttt{rolling.py}, \texttt{reduced.py} y \texttt{descent.py} se testean
primero solamente con \texttt{stadium\_2.json}.
\end{warningbox}

\section{Regla principal de trabajo}

Cada paso debe producir codigo ejecutable y una salida verificable.

Formato obligatorio de reporte:
\begin{verbatim}
Modulo:
Archivo:
Funcion:
Comando:
Salida esperada:
Salida obtenida:
Estado: OK / ERROR
\end{verbatim}

Si el estado es \texttt{ERROR}, no se avanza.

\begin{warningbox}
No se debe avanzar al siguiente modulo si el modulo actual no produce la salida
esperada. La meta no es escribir mucho codigo. La meta es mantener una cadena de
pasos verificados.
\end{warningbox}

\section{Tolerancia numerica}

En todos los tests numericos usamos:
\[
\varepsilon_{\rm test}=10^{-9}.
\]
Una igualdad numerica \(a=b\) significa:
\[
|a-b|<\varepsilon_{\rm test}.
\]
Una igualdad vectorial o matricial \(M=0\) significa:
\[
\|M\|<\varepsilon_{\rm test}.
\]

\begin{warningbox}
No cambiar la tolerancia para hacer pasar un test. Si un test falla con
tolerancia \(10^{-9}\), se corrige el modulo correspondiente.
\end{warningbox}

\section{Regla de errores claros}

Toda funcion debe fallar de manera clara. Si recibe un label, disco, pieza,
segmento, arco, contacto, componente o cruce inexistente, debe levantar
\texttt{ValueError}.

Ejemplos:
\begin{verbatim}
ValueError("Unknown label: alpha")
ValueError("Unknown disk id: 3")
ValueError("Unknown piece: s7")
ValueError("Invalid contact: [1, 5]")
ValueError("Segment endpoint does not exist")
ValueError("Unknown crossing: X3")
ValueError("Component uses unknown piece: a12")
\end{verbatim}

\begin{checkpoint}
Si un dato de entrada es invalido, el programa debe decir donde esta el problema.
No debe fallar despues con un error oscuro.
\end{checkpoint}

\section{Arquitectura de la primera version}

La primera version debe tener esta estructura:

\begin{verbatim}
csdraw_v03/
  csdraw/
    __init__.py
    cs_geometry.py
    length.py
    rolling.py
    reduced.py
    descent.py
    knot_diagram.py
    crossings.py
    variables.py
    validation.py
    optimise.py
    io_json.py
    render_svg.py
  examples/
    stadium_2.json
    figure8_4_1_manual.json
  main.py
  README.md
\end{verbatim}

Los archivos del motor cs son:
\begin{verbatim}
cs_geometry.py
length.py
rolling.py
reduced.py
descent.py
io_json.py
render_svg.py
\end{verbatim}

Los archivos para diagramas de nudos son:
\begin{verbatim}
knot_diagram.py
crossings.py
validation.py
\end{verbatim}

Los archivos para optimizacion son:
\begin{verbatim}
variables.py
optimise.py
\end{verbatim}

\section{Pipeline completo}

El software debe seguir esta flecha:
\[
\boxed{
\text{input}
\to
\text{puntos}
\to
\text{longitud}
\to
\text{SVG}
\to
\text{componentes}
\to
\text{cruces}
\to
4_1
\to
\text{rolling}
\to
\text{reducido}
\to
\text{descenso}
\to
\text{variables}
\to
\text{validacion}
\to
\text{optimizacion}.
}
\]

En la primera semana solo se exige:
\[
\boxed{
\text{stadium\_2 funciona}
\quad\text{y}\quad
\text{figure8\_4\_1\_manual funciona como diagrama con }4\text{ cruces}.
}
\]

\section{No usamos codigo de Gauss en la primera version}

No necesitamos codigo de Gauss para el primer test. La forma mas facil para
Tomas es usar una declaracion explicita de componentes y una tabla explicita de
cruces.

El archivo del nudo \(4_1\) debe almacenar directamente:

\begin{verbatim}
components:
  [
    ["s1", "a1", "s2", "a2", "..."]
  ]

crossings:
  id
  over_piece
  under_piece
  point
\end{verbatim}

Esto quiere decir: el software no intentara reconocer automaticamente el tipo de
nudo desde un codigo de Gauss. En la primera version, el archivo
\texttt{figure8\_4\_1\_manual.json} declara que el dibujo proviene de un
diagrama conocido del nudo \(4_1\), y el programa verifica que la estructura sea
internamente consistente.

\begin{warningbox}
No implementar reconocimiento de nudos en la primera version. No implementar
codigo de Gauss en la primera version. Primero necesitamos medir y validar un
diagrama cs con una componente declarada y cuatro cruces declarados.
\end{warningbox}

\section{Datos cs que usa el programa}

Un diagrama cs se almacena como:

\begin{verbatim}
disks:
  id
  center

labels:
  name
  disk
  theta

segments:
  id
  start
  end

arcs:
  id
  start
  end
  disk
  orientation

contacts:
  [i, j]
\end{verbatim}

La formula central es:
\[
p_\alpha
=
c_{k(\alpha)}
+
(\cos\theta_\alpha,\sin\theta_\alpha).
\]

La longitud es:
\[
\Len
=
\sum_{s:\alpha\to\beta}\|p_\beta-p_\alpha\|
+
\sum_{a:\alpha\to\beta}\Delta_a(\theta_\alpha,\theta_\beta).
\]

\begin{warningbox}
Cada segmento y cada arco debe tener un identificador unico. Ejemplos:
\texttt{s1}, \texttt{s2}, \texttt{a1}, \texttt{a2}. Estos identificadores se
usan para declarar componentes y cruces.
\end{warningbox}

\begin{warningbox}
En esta version, una ``pieza'' significa un segmento o un arco. Por tanto,
\texttt{piece\_ids(diagram)} debe construir la lista de piezas como la union de
los ids en \texttt{segments} y los ids en \texttt{arcs}. No hay una tercera lista
\texttt{pieces} en el JSON.
\end{warningbox}

\section{Datos adicionales para un diagrama de nudo}

Un diagrama de nudo necesita, ademas de los datos cs, una lista de componentes y
una lista de cruces.

La lista de componentes declara el recorrido de la curva:

\begin{verbatim}
components:
  [
    ["s1", "a1", "s2", "a2", "s3", "a3"]
  ]
\end{verbatim}

Para un nudo de una componente, \texttt{components} debe tener longitud \(1\).

\begin{warningbox}
Aunque \texttt{stadium\_2.json} no es un nudo, debe incluir tambien el campo
\texttt{components}. Esto permite testear \texttt{knot\_diagram.py} en un caso
simple antes de usar \texttt{figure8\_4\_1\_manual.json}.
\end{warningbox}

\begin{warningbox}
En esta primera version, las componentes se validan usando el recorrido declarado
de piezas cs. No se reconstruyen a partir de los cruces. Los cruces sirven para
over/under y dibujo, no para construir la componente.
\end{warningbox}

La lista de cruces declara over/under:

\begin{verbatim}
crossings:
  id
  over_piece
  under_piece
  point
\end{verbatim}

Ejemplo esquematico:
\begin{verbatim}
"crossings": [
  {
    "id": "X1",
    "over_piece": "s3",
    "under_piece": "s7",
    "point": [0.0, 1.0]
  }
]
\end{verbatim}

Aqui \texttt{over\_piece} y \texttt{under\_piece} son identificadores de piezas
cs. En la primera version, los cruces se declaran manualmente. El programa solo
verifica consistencia interna.

\begin{warningbox}
En la primera version, el campo \texttt{point} de un cruce es solo una marca
visual del cruce. No se usa para calcular la longitud ni para construir la
componente. Tampoco se exige todavia verificar geometricamente que las dos
piezas se intersecten exactamente en ese punto. Esa validacion se agrega
despues.
\end{warningbox}

\section{Micro-modulo previo: construir los archivos JSON de test}

\subsection{Archivo \texttt{stadium\_2.json}}

El archivo \texttt{stadium\_2.json} debe ser un test cs simple, pero debe tener
los mismos campos estructurales que un diagrama de nudo cuando esto sea posible.

Debe incluir:

\begin{verbatim}
"type": "cs_loop",
"name": "stadium_2",
"components": [["s1", "a1", "s2", "a2"]],
"crossings": []
\end{verbatim}

\begin{checkpoint}
El estadio no tiene cruces, pero debe tener \texttt{crossings: []}. Asi el
codigo de cruces no necesita tratar el estadio como caso especial.
\end{checkpoint}

\subsection{Archivo \texttt{figure8\_4\_1\_manual.json}}

El archivo \texttt{figure8\_4\_1\_manual.json} debe tener discos, labels,
segmentos, arcos, components y crossings.

No se exige todavia que el dibujo sea perfecto. Se exige que el JSON sea
legible, cerrado y tenga cuatro cruces declarados.

Debe incluir:

\begin{verbatim}
"type": "knot_diagram",
"name": "figure8_4_1_manual",
"knot_label": "4_1"
\end{verbatim}

Debe incluir exactamente una componente:

\begin{verbatim}
"components": [
  ["...", "...", "..."]
]
\end{verbatim}

Debe incluir exactamente cuatro cruces:

\begin{verbatim}
"crossings": [
  {"id": "X1", "over_piece": "...", "under_piece": "...", "point": [...]},
  {"id": "X2", "over_piece": "...", "under_piece": "...", "point": [...]},
  {"id": "X3", "over_piece": "...", "under_piece": "...", "point": [...]},
  {"id": "X4", "over_piece": "...", "under_piece": "...", "point": [...]}
]
\end{verbatim}

\begin{checkpoint}
Antes de programar \texttt{crossings.py}, abrir el JSON y verificar visualmente
que existen exactamente cuatro ids de cruces: \texttt{X1}, \texttt{X2},
\texttt{X3}, \texttt{X4}.
\end{checkpoint}

\section{Que significa verificar \(4_1\) en esta version}

El programa no demuestra automaticamente que el nudo sea \(4_1\). El archivo
\texttt{figure8\_4\_1\_manual.json} debe construirse a partir de un diagrama
conocido de \(4_1\), incluyendo sus piezas, su recorrido de componente y sus
cruces over/under. El programa verifica que esos datos sean internamente
consistentes.

El test de \(4_1\) debe verificar:

\begin{enumerate}[label=\textbf{\arabic*.}]
\item el archivo tiene \texttt{"type": "knot\_diagram"};
\item el archivo tiene \texttt{"knot\_label": "4\_1"};
\item hay exactamente cuatro cruces declarados;
\item cada cruce tiene id unico;
\item cada cruce tiene una pieza over y una pieza under;
\item las piezas over/under existen;
\item la pieza over y la pieza under no son la misma;
\item hay exactamente una componente declarada;
\item la componente usa solo piezas existentes;
\item cada pieza de \texttt{segments} y \texttt{arcs} aparece exactamente una vez
en alguna componente;
\item la componente es cerrada;
\item la longitud cs se puede calcular;
\item se puede exportar SVG.
\end{enumerate}

\begin{checkpoint}
Para \texttt{figure8\_4\_1\_manual.json}, el programa debe reportar:
\begin{verbatim}
knot label: 4_1
number of crossings: 4
number of components: 1
component pieces valid: True
all pieces used exactly once: True
component closed: True
crossing ids valid: True
crossing pieces valid: True
over/under valid: True
crossings valid: True
length valid: True
\end{verbatim}
\end{checkpoint}

\begin{checkpoint}
Para \texttt{figure8\_4\_1\_manual.json}, la longitud debe ser finita y positiva:
\[
0<\Len<+\infty.
\]
No se exige un valor numerico especifico en la primera version.
\end{checkpoint}

\begin{checkpoint}
Para \texttt{figure8\_4\_1\_manual.json}, si el JSON no cambia, la longitud debe
ser reproducible: dos ejecuciones consecutivas deben dar el mismo valor
numerico.
\end{checkpoint}

\section{Test obligatorio 1: \texttt{stadium\_2.json}}

\subsection{Rol del test}

\texttt{stadium\_2.json} es el primer test obligatorio. No es un nudo. Es el
test unitario del motor geometrico cs.

Verifica:

\begin{enumerate}[label=\textbf{\arabic*.}]
\item lectura JSON;
\item discos y labels;
\item puntos \(p_\alpha\);
\item segmentos;
\item arcos;
\item longitud;
\item SVG;
\item componentes declaradas;
\item cruces vacios;
\item contactos;
\item matriz \(A(c)\);
\item kernel de \(A(c)\);
\item gradiente reducido;
\item test de descenso.
\end{enumerate}

\subsection{Datos}

\[
c_1=(0,0),
\qquad
c_2=(2,0).
\]

Hay dos segmentos de longitud \(2\) y dos semicirculos de longitud \(\pi\).
Por tanto:
\[
\Len=2+2+\pi+\pi=4+2\pi.
\]

Valor numerico esperado:
\[
\Len=10.283185307179586.
\]

\subsection{Comando}

\begin{verbatim}
python main.py examples/stadium_2.json
\end{verbatim}

\subsection{Salida esperada}

\begin{verbatim}
Diagram: stadium_2
Type: cs_loop
Number of disks: 2
Number of labels: 4
Number of segments: 2
Number of arcs: 2
Number of components: 1
Number of crossings: 0
Total length: 10.283185307179586
Rolling matrix A:
[[-1.,  0.,  1., -0.]]
rank A: 1
dim Roll: 3
shape K: (4, 3)
norm A K: < 1e-9
g_red: [-2, 0, 2, 0]
stationary: True
has descent: False
SVG written to: examples/stadium_2.svg
\end{verbatim}

\begin{checkpoint}
Si \texttt{stadium\_2.json} falla, no se prueba \(4_1\). Primero se corrige el
motor geometrico.
\end{checkpoint}

\section{Test obligatorio 2: \texttt{figure8\_4\_1\_manual.json}}

\subsection{Rol del test}

\texttt{figure8\_4\_1\_manual.json} es el primer test real de diagrama de nudo.
No prueba reconocimiento automatico. Prueba que el software puede leer, medir,
dibujar y validar un diagrama cs de nudo con cruces.

\subsection{Datos esperados}

El archivo debe tener:

\begin{verbatim}
"type": "knot_diagram"
"knot_label": "4_1"
\end{verbatim}

Debe tener exactamente una componente declarada:

\begin{verbatim}
components:
  [
    ["...", "...", "..."]
  ]
\end{verbatim}

Debe tener cuatro cruces declarados:

\begin{verbatim}
crossings:
  [
    {"id": "X1", "over_piece": "...", "under_piece": "...", "point": [...]},
    {"id": "X2", "over_piece": "...", "under_piece": "...", "point": [...]},
    {"id": "X3", "over_piece": "...", "under_piece": "...", "point": [...]},
    {"id": "X4", "over_piece": "...", "under_piece": "...", "point": [...]}
  ]
\end{verbatim}

\subsection{Comando}

\begin{verbatim}
python main.py examples/figure8_4_1_manual.json
\end{verbatim}

\subsection{Salida esperada}

\begin{verbatim}
Diagram: figure8_4_1_manual
Type: knot_diagram
Knot label: 4_1
Number of crossings: 4
Number of components: 1
Component pieces valid: True
All pieces used exactly once: True
Component closed: True
Crossing ids valid: True
Crossing pieces valid: True
Over/under valid: True
Crossings valid: True
One component: True
Total length:
SVG written to: examples/figure8_4_1_manual.svg
\end{verbatim}

\begin{checkpoint}
Si \texttt{figure8\_4\_1\_manual.json} falla, no pasar a optimizacion. Primero
corregir componentes, cruces, piezas o JSON.
\end{checkpoint}

\section{Modulo 1: \texttt{cs\_geometry.py}}

\subsection{Responsabilidad}

Este modulo define la geometria cs:

\begin{verbatim}
Disk
Label
SegmentPiece
ArcPiece
CSRealization
\end{verbatim}

Su funcion mas importante es:

\begin{verbatim}
point(label_name)
\end{verbatim}

Debe calcular:
\[
p_\alpha=c_{k(\alpha)}+(\cos\theta_\alpha,\sin\theta_\alpha).
\]

\subsection{Test}

Para el estadio:
\[
p_\alpha=(0,1),
\quad
p_\beta=(0,-1),
\quad
p_\gamma=(2,1),
\quad
p_\delta=(2,-1).
\]

\begin{checkpoint}
Si estos cuatro puntos no salen bien, no revisar longitud. Primero corregir
\texttt{cs\_geometry.py}.
\end{checkpoint}

\section{Modulo 2: \texttt{length.py}}

\subsection{Responsabilidad}

Este modulo calcula:

\begin{verbatim}
segment_length(cs, s)
arc_length(cs, a)
total_length(cs)
length_report(cs)
\end{verbatim}

Formula de segmento:
\[
\ell_s=\|p_\beta-p_\alpha\|.
\]

Formula de arco:
\[
\ell_a=\Delta_a(\theta_\alpha,\theta_\beta).
\]

\subsection{Test}

Para el estadio:
\[
\ell_{\alpha\gamma}=2,
\qquad
\ell_{\delta\beta}=2,
\qquad
\ell_{\gamma\delta}=\pi,
\qquad
\ell_{\beta\alpha}=\pi.
\]

Total:
\[
\Len=10.283185307179586.
\]

\begin{checkpoint}
Si la longitud total falla, corregir \texttt{length.py}. No tocar los modulos de
componentes ni cruces.
\end{checkpoint}

\section{Modulo 3: \texttt{knot\_diagram.py}}

\subsection{Responsabilidad}

Este modulo valida la estructura de componentes y permite tratar el archivo como
diagrama de nudo.

No reconoce el tipo de nudo. Solo verifica estructura minima.

\subsection{Funciones}

\begin{verbatim}
validate_knot_metadata(diagram)
piece_ids(diagram)
piece_graph(diagram)
validate_components(diagram)
all_pieces_used_once(diagram)
component_count(diagram)
component_closed(diagram, component)
is_one_component(diagram)
knot_diagram_report(diagram)
\end{verbatim}

\subsection{Micro-modulo K0: validar metadata de nudo}

Funcion:
\begin{verbatim}
validate_knot_metadata(diagram)
\end{verbatim}

Debe verificar que, para \texttt{figure8\_4\_1\_manual.json},
\begin{verbatim}
type == "knot_diagram"
knot_label == "4_1"
\end{verbatim}

\begin{checkpoint}
Si el campo \texttt{type} o \texttt{knot\_label} falta, el programa debe levantar
\texttt{ValueError} con un mensaje claro.
\end{checkpoint}

\subsection{Micro-modulo K1: lista de piezas}

Funcion:
\begin{verbatim}
piece_ids(diagram)
\end{verbatim}

Debe devolver todos los identificadores de piezas cs. Una pieza es un segmento o
un arco. Por ejemplo:

\begin{verbatim}
["s1", "s2", "a1", "a2"]
\end{verbatim}

\begin{checkpoint}
Cada segmento y cada arco debe tener un identificador unico.
\end{checkpoint}

\subsection{Micro-modulo K2: validar componentes}

Funcion:
\begin{verbatim}
validate_components(diagram)
\end{verbatim}

Debe verificar:

\begin{enumerate}[label=\textbf{\arabic*.}]
\item existe el campo \texttt{components};
\item cada componente es una lista de piezas;
\item cada pieza mencionada existe;
\item ninguna componente esta vacia;
\item cada pieza de \texttt{segments} y \texttt{arcs} aparece exactamente una vez
en alguna componente.
\end{enumerate}

\begin{checkpoint}
Para \texttt{figure8\_4\_1\_manual.json}, debe devolver:
\begin{verbatim}
component pieces valid: True
all pieces used exactly once: True
\end{verbatim}
\end{checkpoint}

\subsection{Micro-modulo K3: componente cerrada}

Funcion:
\begin{verbatim}
component_closed(diagram, component)
\end{verbatim}

Debe verificar que las piezas se puedan recorrer consecutivamente.

Para cualquier pieza \(P\), sea \(\operatorname{start}(P)\) su label inicial y
\(\operatorname{end}(P)\) su label final. Si \(P\) es un segmento, estos campos
vienen de \texttt{segments}. Si \(P\) es un arco, estos campos vienen de
\texttt{arcs}. La funcion \texttt{component\_closed} no distingue entre
segmentos y arcos: solo usa \texttt{start} y \texttt{end}.

Si una pieza \(P_i\) termina en el label \(\beta\), la siguiente pieza
\(P_{i+1}\) debe empezar en \(\beta\). Ademas, la ultima pieza debe terminar en
el label donde empieza la primera.

\begin{checkpoint}
Para el estadio y para \(4_1\), debe devolver:
\begin{verbatim}
component closed: True
\end{verbatim}
\end{checkpoint}

\subsection{Micro-modulo K4: contar componentes}

Funcion:
\begin{verbatim}
component_count(diagram)
\end{verbatim}

Para el estadio:
\[
\#\text{componentes}=1.
\]

Para \(4_1\):
\[
\#\text{componentes}=1.
\]

\begin{checkpoint}
El test \(4_1\) debe decir:
\begin{verbatim}
one component: True
\end{verbatim}
\end{checkpoint}

\section{Modulo 4: \texttt{crossings.py}}

\subsection{Responsabilidad}

Este modulo valida los cruces declarados manualmente.

No debe reconocer nudos. No debe usar codigo de Gauss. Solo verifica que la
tabla de cruces sea consistente.

\subsection{Funciones}

\begin{verbatim}
crossing_count(diagram)
validate_crossing_ids(diagram)
validate_crossing_pieces(diagram)
validate_over_under(diagram)
crossing_report(diagram)
\end{verbatim}

\subsection{Micro-modulo C1: contar cruces}

Funcion:
\begin{verbatim}
crossing_count(diagram)
\end{verbatim}

Para el estadio:
\[
\#\text{cruces}=0.
\]

Para \(4_1\):
\[
\#\text{cruces}=4.
\]

\begin{checkpoint}
El reporte debe decir:
\begin{verbatim}
stadium_2: number of crossings = 0
figure8_4_1_manual: number of crossings = 4
\end{verbatim}
\end{checkpoint}

\subsection{Micro-modulo C2: ids de cruces}

Funcion:
\begin{verbatim}
validate_crossing_ids(diagram)
\end{verbatim}

Debe verificar que cada cruce tenga un id y que los ids no se repitan.

\begin{checkpoint}
Para \(4_1\), los ids esperados son \texttt{X1}, \texttt{X2}, \texttt{X3},
\texttt{X4}, sin repeticiones.
\end{checkpoint}

\subsection{Micro-modulo C3: validar piezas de cruce}

Funcion:
\begin{verbatim}
validate_crossing_pieces(diagram)
\end{verbatim}

Debe verificar que cada \texttt{over\_piece} y cada \texttt{under\_piece}
existan como piezas cs.

\begin{checkpoint}
Si un cruce menciona una pieza inexistente, debe levantar \texttt{ValueError}.
\end{checkpoint}

\subsection{Micro-modulo C4: validar over/under}

Funcion:
\begin{verbatim}
validate_over_under(diagram)
\end{verbatim}

Debe verificar:

\begin{enumerate}[label=\textbf{\arabic*.}]
\item cada cruce tiene una pieza over;
\item cada cruce tiene una pieza under;
\item la pieza over y la pieza under no son la misma;
\item el punto del cruce tiene dos coordenadas reales.
\end{enumerate}

\begin{checkpoint}
Para \texttt{figure8\_4\_1\_manual.json}, debe devolver \texttt{True}.
\end{checkpoint}

\subsection{Micro-modulo C5: reporte de cruces}

Funcion:
\begin{verbatim}
crossing_report(diagram)
\end{verbatim}

Salida esperada para \(4_1\):
\begin{verbatim}
number of crossings: 4
crossing ids valid: True
crossing pieces valid: True
over/under valid: True
crossings valid: True
\end{verbatim}

\begin{taskbox}
Actualizar \texttt{main.py} para imprimir \texttt{crossing\_report} si el JSON
contiene una lista \texttt{crossings}.
\end{taskbox}

\begin{checkpoint}
No pasar a rolling, reduced ni descent para \(4_1\) hasta que
\texttt{crossing\_report} funcione para \(4_1\).
\end{checkpoint}

\section{Modulo 5: reporte de diagrama de nudo}

\subsection{Responsabilidad}

El reporte de diagrama de nudo junta la informacion de metadata, componentes,
cruces, longitud y SVG.

\subsection{Funcion}

\begin{verbatim}
knot_diagram_report(diagram)
\end{verbatim}

Salida esperada para \(4_1\):
\begin{verbatim}
Diagram: figure8_4_1_manual
Type: knot_diagram
Knot label: 4_1
Number of crossings: 4
Number of components: 1
Component pieces valid: True
All pieces used exactly once: True
Component closed: True
Crossings valid: True
One component: True
Total length:
SVG written to: examples/figure8_4_1_manual.svg
\end{verbatim}

\begin{taskbox}
Actualizar \texttt{main.py} para llamar \texttt{knot\_diagram\_report} cuando el
JSON tenga:
\begin{verbatim}
"type": "knot_diagram"
\end{verbatim}
\end{taskbox}

\begin{checkpoint}
Este es el primer test real del proyecto como software para diagramas cs de
nudos.
\end{checkpoint}

\section{Modulo 6: \texttt{rolling.py}}

\subsection{Responsabilidad}

Este modulo calcula la matriz de rolling \(A(c)\).

Para un contacto \(\{i,j\}\), se define:
\[
u_{ij}=\frac{c_i-c_j}{\|c_i-c_j\|}.
\]

La restriccion linealizada es:
\[
\langle u_{ij},\delta c_i-\delta c_j\rangle=0.
\]

Entonces \(A(c)\) tiene una fila por contacto.

\subsection{Funciones}

\begin{verbatim}
build_A(cs, contacts)
contact_residuals(cs, contacts)
validate_contacts(cs, contacts)
kernel_basis(A)
rolling_report(cs, contacts)
\end{verbatim}

\subsection{Test para estadio}

Para el estadio:
\[
A(c)=
\begin{pmatrix}
-1&0&1&0
\end{pmatrix},
\qquad
\rank A=1,
\qquad
\dim\Roll(c)=3.
\]

Ademas, si \(K\) genera \(\Ker A(c)\), debe cumplirse:
\[
K\in\R^{4\times 3},
\qquad
\|AK\|<10^{-9}.
\]

\begin{checkpoint}
Si \(A(c)\) o \(K\) fallan para el estadio, no pasar a \texttt{reduced.py}.
\end{checkpoint}

\section{Modulo 7: \texttt{reduced.py}}

\subsection{Responsabilidad}

Este modulo calcula el gradiente reducido \(g_{\rm red}\) y el test reducido de
estacionariedad.

No usa gauge. No usa \(\omega\). Solo usa centros, segmentos y \(A(c)\).

\subsection{Funciones}

\begin{verbatim}
assemble_g_red(cs)
reduced_kernel_residual(cs, contacts)
row_space_residual(cs, contacts)
reduced_stationarity_report(cs, contacts)
\end{verbatim}

\subsection{Test para estadio}

Para el estadio:
\[
g_{\rm red}=(-2,0,2,0)^T.
\]

Tambien:
\[
\|g_{\rm red}^TK\|<10^{-9}.
\]

Y:
\[
g_{\rm red}=2A(c)^T.
\]

Salida esperada:
\begin{verbatim}
g_red: [-2, 0, 2, 0]
shape K: (4, 3)
norm A K: < 1e-9
rank A: 1
dim Roll: 3
kernel residual: < 1e-9
row-space residual: < 1e-9
stationary: True
\end{verbatim}

\begin{checkpoint}
No pasar a \texttt{descent.py} hasta que el estadio sea estacionario.
\end{checkpoint}

\section{Modulo 8: \texttt{descent.py}}

\subsection{Responsabilidad}

Este modulo encuentra una direccion de descenso cuando el diagrama no es
estacionario.

Si \(K\) genera \(\Roll(c)\), toda direccion admisible lineal tiene la forma:
\[
\delta c=Ky.
\]

Si \(g_{\rm red}^TK\neq0\), una direccion de descenso es:
\[
\delta c=-KK^Tg_{\rm red}.
\]

\subsection{Funciones}

\begin{verbatim}
find_descent_direction(cs, contacts)
descent_report(cs, contacts)
\end{verbatim}

\subsection{Test para estadio}

Como el estadio es estacionario:
\[
\|\text{gradiente proyectado}\|<10^{-9}.
\]

Salida esperada:
\begin{verbatim}
projected gradient norm: < 1e-9
has descent: False
\end{verbatim}

\begin{checkpoint}
No pasar a variables u optimizacion si este test falla.
\end{checkpoint}

\section{Modulo 9: \texttt{variables.py}}

\subsection{Responsabilidad}

Este modulo empaqueta y desempaqueta variables para optimizacion no lineal.

Variables:
\[
x=(c_{1x},c_{1y},\ldots,c_{Nx},c_{Ny},
\theta_{\alpha_1},\ldots,\theta_{\alpha_M}).
\]

\subsection{Funciones}

\begin{verbatim}
pack_variables(cs)
unpack_variables(x, template)
variable_report(cs)
\end{verbatim}

\begin{warningbox}
El orden de las etiquetas debe ser el mismo en \texttt{pack\_variables},
\texttt{unpack\_variables}, \texttt{assemble\_g\_red} y todos los reportes.
\end{warningbox}

\subsection{Test para estadio}

Para el estadio:
\[
N=2,\qquad M=4,\qquad x\in\R^8.
\]

Salida esperada:
\begin{verbatim}
number of disks: 2
number of labels: 4
number of variables: 8
pack/unpack length error: < 1e-9
\end{verbatim}

\begin{checkpoint}
No pasar a \texttt{optimise.py} hasta que pack/unpack conserve la longitud.
\end{checkpoint}

\section{Modulo 10: \texttt{validation.py}}

\subsection{Responsabilidad}

Este modulo valida que una configuracion siga siendo admisible.

\subsection{Funciones}

\begin{verbatim}
check_contacts(cs, contacts)
check_disk_clearance(cs, contacts)
minimum_clearance(cs, contacts)
validation_report(cs, contacts)
\end{verbatim}

\subsection{Tests}

Para el estadio:
\[
\|c_1-c_2\|-2=0.
\]

Para \(4_1\), el reporte debe verificar:

\begin{enumerate}[label=\textbf{\arabic*.}]
\item metadata valida;
\item cruces validos;
\item una componente;
\item todas las piezas usadas exactamente una vez;
\item componente cerrada;
\item contactos validos, si hay contactos prescritos;
\item no solapamiento de discos;
\item longitud computable.
\end{enumerate}

Salida esperada para \(4_1\):
\begin{verbatim}
metadata valid: True
crossings valid: True
one component: True
all pieces used exactly once: True
component closed: True
contacts valid: True
disk clearance valid: True
length valid: True
valid: True
\end{verbatim}

\begin{checkpoint}
No pasar a optimizacion si \texttt{validation\_report} no dice
\texttt{valid: True}.
\end{checkpoint}

\section{Modulo 11: \texttt{optimise.py}}

\subsection{Responsabilidad}

Este modulo minimiza la longitud para un tipo cs fijo.

No se debe escribir antes de que funcionen:

\begin{verbatim}
length.py
knot_diagram.py
crossings.py
rolling.py
reduced.py
descent.py
variables.py
validation.py
\end{verbatim}

\begin{warningbox}
En la primera semana no se exige obtener un minimo. Solo se exige que
\texttt{objective\_length}, \texttt{contact\_constraints} y
\texttt{clearance\_constraints} funcionen y pasen sus tests.
\end{warningbox}

\subsection{Funciones}

\begin{verbatim}
objective_length(x, template)
contact_constraints(x, template, contacts)
clearance_constraints(x, template, contacts)
minimise_length(template, contacts)
optimisation_report(result, template, contacts)
\end{verbatim}

\subsection{Micro-modulo O1: objetivo de longitud}

Funcion:
\begin{verbatim}
objective_length(x, template)
\end{verbatim}

Debe desempaquetar \(x\), construir la realizacion cs y devolver \(\Len\).

Test para estadio:
\begin{verbatim}
x = pack_variables(cs)
objective_length(x, cs) = 10.283185307179586
\end{verbatim}

\begin{checkpoint}
Si este test falla, no implementar restricciones.
\end{checkpoint}

\subsection{Micro-modulo O2: restricciones de contacto}

Para cada contacto:
\[
h_{ij}(x)=\|c_i-c_j\|-2.
\]

Debe imponerse:
\[
h_{ij}(x)=0.
\]

\subsection{Micro-modulo O3: restricciones de clearance}

Para cada no contacto:
\[
q_{ij}(x)=\|c_i-c_j\|-2.
\]

Debe imponerse:
\[
q_{ij}(x)\geq0.
\]

\begin{warningbox}
Contactos prescritos son igualdades. No contactos son desigualdades. No mezclar.
\end{warningbox}

\subsection{Micro-modulo O4: primer minimizador}

Usar inicialmente:
\begin{verbatim}
scipy.optimize.minimize(method="SLSQP")
\end{verbatim}

\begin{checkpoint}
No correr SLSQP hasta que \texttt{objective\_length},
\texttt{contact\_constraints} y \texttt{clearance\_constraints} funcionen.
\end{checkpoint}

\section{Plan de trabajo de 7 dias}

\subsection*{Dia 1: motor cs basico}

Correr:
\begin{verbatim}
python main.py examples/stadium_2.json
\end{verbatim}

Entregar:
\begin{verbatim}
Total length = 10.283185307179586
SVG generado
\end{verbatim}

\subsection*{Dia 2: componentes declaradas}

Implementar:
\begin{verbatim}
knot_diagram.py
\end{verbatim}

Entregar:
\begin{verbatim}
stadium_2: number of components = 1
figure8_4_1_manual: number of components = 1
component pieces valid: True
all pieces used exactly once: True
component closed: True
\end{verbatim}

\subsection*{Dia 3: cruces manuales}

Implementar:
\begin{verbatim}
crossings.py
\end{verbatim}

Entregar:
\begin{verbatim}
stadium_2: number of crossings = 0
figure8_4_1_manual: number of crossings = 4
crossing ids valid: True
crossings valid: True
\end{verbatim}

\subsection*{Dia 4: primer reporte de nudo}

Actualizar:
\begin{verbatim}
main.py
\end{verbatim}

Entregar:
\begin{verbatim}
Diagram: figure8_4_1_manual
Knot label: 4_1
Number of crossings: 4
Number of components: 1
All pieces used exactly once: True
Component closed: True
One component: True
Total length:
SVG written
\end{verbatim}

\subsection*{Dia 5: rolling, reducido y descenso}

Implementar o verificar:
\begin{verbatim}
rolling.py
reduced.py
descent.py
\end{verbatim}

Entregar para estadio:
\begin{verbatim}
A matrix
rank A
dim Roll
shape K
norm A K < 1e-9
g_red = [-2, 0, 2, 0]
stationary: True
has descent: False
\end{verbatim}

\begin{warningbox}
En el Dia 5, rolling, reduced y descent se testean primero solo en
\texttt{stadium\_2.json}. No se exige todavia que el gradiente reducido sea
estacionario para \(4_1\).
\end{warningbox}

\subsection*{Dia 6: variables y validacion}

Implementar:
\begin{verbatim}
variables.py
validation.py
\end{verbatim}

Entregar:
\begin{verbatim}
pack/unpack conserva longitud
validation_report pasa para stadium_2
validation_report pasa para figure8_4_1_manual
\end{verbatim}

\subsection*{Dia 7: preparar optimizacion}

Implementar:
\begin{verbatim}
optimise.py
\end{verbatim}

Solo hasta:
\begin{verbatim}
objective_length
contact_constraints
clearance_constraints
\end{verbatim}

No correr SLSQP todavia.

\section{Regla final}

Este proyecto es para construir software de diagramas cs de nudos. Por eso hay
dos tests obligatorios:

\[
\boxed{
\texttt{stadium\_2.json}
=
\text{test geometrico minimo}
}
\]

\[
\boxed{
\texttt{figure8\_4\_1\_manual.json}
=
\text{primer test real de diagrama de nudo}
}
\]

La cadena de verdad es:
\[
\boxed{
\text{cs engine}
\to
\text{components}
\to
\text{crossings}
\to
4_1
\to
\text{rolling}
\to
\text{reduced}
\to
\text{descent}
\to
\text{validation}
\to
\text{optimisation}.
}
\]

Si una flecha falla, se corrige esa flecha antes de avanzar.

\end{document}